\section{Threats to Validity}\label{sec:threats}

The contribution of this article is a reference architecture, validated by instantiation rather than by a controlled empirical study, and the honest assessment of what that instantiation can and cannot establish is itself part of the contribution. This section organizes the threats along the conventional four axes---construct, internal, external, and conclusion validity---and, for each threat, names the design mechanism intended to contain it. The recurring theme is that several of the architecture's most distinctive features---provenance- and assurance-tagged obligations, the independent-authority constraint, behaviour-level ratification, and tiered assurance---exist precisely because the threats below are real and were anticipated, rather than discovered after the fact. Where a threat cannot be eliminated, the design aims to make it \emph{visible and auditable} rather than to claim it away.

\subsection{Construct Validity}
\label{subsec:threats-construct}

Construct validity concerns whether the artefacts the pipeline produces actually measure what the lifecycle requires of them---namely, a faithful, machine-checkable encoding of intent. Three constructs are at risk.

The most fundamental is that \emph{inferred specifications skew weak}. Heuristic AST-based inference, and dynamic invariant detection of the Daikon lineage~\cite{ernst2007daikon} more generally, describe what the code \emph{does}, not what it \emph{should} do; a draft contract recovered from an implementation that silently truncates on overflow will faithfully record the truncation as intended behaviour. A weak contract is a poor oracle: code can satisfy it while still violating the human's real intent. The architecture does not pretend that inference closes this gap. Instead it treats the inferred contract as an explicitly \emph{code-inferred}, \emph{unconfirmed} provenance class that \emph{drives} the human dialogue rather than gating it, and it strengthens weak drafts through three channels: inherited default contracts that impose organisation-level policy (for example, public methods reject \texttt{null} unless stated) the elaboration only needs to override on departure; adversarial elaboration, the elaboration-layer dual of counterexample-guided inductive synthesis (CEGIS)~\cite{solarlezama2008sketching}, in which a second agent searches for scenarios consistent with the natural-language intent but admitted by the proposed contract, flagging weakenings such as ``sort'' degraded to ``is a permutation''; and vacuity and non-triviality checks that repurpose the author's own ``ensures true'' punt detection as a weakening signal. The residual threat---that no automated check substitutes for intent the human never articulated---is mitigated only by ratification, which is real labor and is itself bounded below.

A dual construct threat is that inferred specifications are sometimes \emph{wrong or over-strong}, producing \emph{false candidate bugs}. An over-fitted clause derived from a single code path may declare a behavioural pattern the method does not in fact guarantee, so that legitimate inputs appear to violate it. Because the brownfield startup phase reports disagreements between code and its own inferred pattern as \emph{candidate} bugs, never as confirmed defects, and surfaces them as the opening agenda of the human dialogue rather than as a gate, an over-strong clause costs a conversation, not a false alarm in production. The provenance tagging makes the source of each clause auditable, so a reviewer can discount a \emph{code-inferred} candidate appropriately.

The third construct threat is that the \emph{natural-language-to-formal translation is imperfect}. Converting tickets and documentation into formal clauses without a human in the loop is an open research problem~\cite{cosler2023nl2spec,endres2024nl2postcond,ma2025specgen}, and a mistranslation injects a defective oracle. The design constrains the blast radius of this threat structurally: any clause derived from a natural-language source without human confirmation retains \emph{ticket-stated} or \emph{doc-stated} provenance and \emph{unconfirmed} assurance, so it never silently becomes a trusted oracle---it enters the ratification dialogue as a question, and only behaviour-level confirmation promotes it.

\subsection{Internal Validity}
\label{subsec:threats-internal}

Internal validity concerns whether the causal mechanism the article claims---that an independent, machine-checkable contract supplies the verification spine the surveyed agentic lifecycles lack---actually holds within the instantiation, rather than being an artefact of how the pieces are wired together.

The dominant internal threat is \emph{circularity}, in two forms. The first is the self-grading hazard endemic to closed-loop LLM systems, in which the same model authors both the implementation and the oracle against which it is checked~\cite{liu2026oracle,sun2024clover}. If contract and code share an author, a model that misreads the intent produces an oracle that excuses its own error, and apparent verification is vacuous. The architecture answers this with the \emph{independent-authority constraint}: the contract the code is verified against must not be authored by the same agent instance or role that writes the code. The elaborator role is maximally faithful and strict, never sees the implementation, and the implementer role is tasked only with satisfying a contract it did not write. Crucially, the verifier is OpenJML with Z3~\cite{cok2014openjml,demoura2008z3}, a trusted symbolic oracle that never grades itself, so the soundness arbiter is structurally outside the LLM loop. Provenance tags make the constraint \emph{auditable}: a verification claim resting on a clause whose provenance is the implementer's own role is detectable and rejectable. The residual threat is that role separation can be subverted by shared context or prompt leakage; the design treats this as a configuration invariant to be checked, not as a property assumed for free.

The second, subtler internal threat is \emph{verifying against the wrong specification}---the case in which the contract is internally consistent, the code provably satisfies it, and yet the contract does not encode the human's intent. A wrong spec verified against is worse than no spec, because it launders an incorrect implementation behind a green proof. This is exactly why behaviour-level ratification and the brownfield triangulation are load-bearing rather than ornamental. The pipeline never asks the human to confirm formal text---which would make the human authority fake, since the human resorted to natural language precisely because they cannot or will not read contracts (the \emph{ratification gap})---but instead renders the contract as concrete, checkable examples (``given $x=-1$ this contract says the result is $0$; is that correct?''), and ratified examples become \emph{pinned oracles} the contract must continue to satisfy. In the brownfield setting, correctness relative to intent is sought not from the code-derived spec alone---verifying code against a code-derived spec is \emph{characterisation}, a faithful-description baseline and toolchain sanity check that is nearly circular and proves nothing about absence of bugs---but from \emph{disagreement} among the three evidence streams: code, tickets, and documentation. A defect surfaces as code violating its own inferred pattern or, more tellingly, code disagreeing with ticket- or doc-derived intent.

A further internal threat is the \emph{noisiness of ticket- and doc-to-code traceability}. Linking a Jira ticket or a Javadoc paragraph to the method it concerns relies on commit links, path heuristics, and information-retrieval techniques that are known to be imperfect~\cite{antoniol2002recovering,borg2014recovering,lin2021tbert}, and comment-code inconsistency is itself well documented~\cite{tan2012tcomment,wen2019codecomment}. Mislinked intent corrupts the disagreement signal on which candidate-bug detection depends. The design contains this threat by retrieval-augmented heuristics whose output is, again, never a trusted oracle: a traced clause is \emph{ticket-stated} and \emph{unconfirmed} until ratified, so a spurious link produces a question for the human, not a silent corruption. The article additionally restricts strong claims to a curated module surface where traceability is credible, rather than asserting corpus-wide link accuracy.

\subsection{External Validity}
\label{subsec:threats-external}

External validity concerns the generality of the demonstration beyond the conditions under which it was produced.

The principal external threat is the \emph{single-corpus} basis of the proof-of-concept. The instantiation is exercised on Apache Commons Lang, chosen because it couples a real public issue tracker (the Apache JIRA \texttt{LANG-\#\#\#\#} project) with rich Javadoc, yielding a genuine three-stream code--tickets--docs dataset of the kind the brownfield phase requires. A single, well-maintained, mature Java library is not representative of the breadth of production software: its API discipline, documentation quality, and ticket hygiene are higher than typical, which flatters traceability and characterisation alike. The article therefore claims \emph{feasibility and realism} of the reference architecture, not a controlled effect size, and explicitly defers validation on external production codebases and an industrial case study---action research under human-research-ethics approval---to future work.

A second external threat is the \emph{decidability and scalability ceiling of verification}, which bounds the class of programs to which the pipeline's strongest assurance applies. Extended static checking discharged by an SMT solver does not always terminate: Z3 times out on multiplicative reasoning, array-theory queries, and rich preconditions, as the author's own prior experiments document, and a custom OpenJML fork was required merely to model recursive \verb|\sum|, \verb|\product|, and \verb|\num_of|. The instantiation consequently targets first the method classes that already verify green---null-safety, array-bounds, simple arithmetic postconditions, and the loop patterns (max/min, linear search, sum) for which the inferrer already produces invariants---and adopts \emph{frontier adoption} for legacy scale, specifying only the module under change, the public API surface, and modules referenced by active tickets while treating the remainder as unspecified under weak inherited defaults. The mitigation for non-termination is \emph{tiered assurance with graceful degradation}: full proof degrades to bounded check, then to test, with the achieved assurance level \emph{recorded} on each obligation and never silently dropped, so a timeout downgrades a claim rather than invalidating the pipeline.

\subsection{Conclusion Validity}
\label{subsec:threats-conclusion}

Conclusion validity concerns whether the inferences drawn from the demonstration are warranted, and is governed here chiefly by \emph{LLM nondeterminism}. The elicitation, contract-drafting, and synthesis steps invoke a stochastic model, so a single run is not a stable measurement: a contract drafted, an ambiguity surfaced, or an implementation synthesized on one invocation may differ on the next, and reported behaviour could reflect a favorable sample. The architecture localizes the consequences of this nondeterminism by routing every consequential output through the deterministic verifier and the deterministic provenance ledger: a nondeterministic implementation is admitted only if it satisfies the contract under OpenJML, and a nondeterministically drafted clause remains \emph{unconfirmed} until ratified into a pinned oracle, so variability degrades to retry-until-verified rather than to silent acceptance of an unsound result. The synthesis--verification loop is capped at $N$ iterations, after which the failing obligation is reported under tiered assurance rather than masked. What this does \emph{not} neutralize is the threat to any quantitative claim: the article accordingly frames its results as an existence demonstration of the loop on a curated corpus, and any future effect-size claim will require repeated runs with reported variance, multiple corpora, and the ethics-approved case study noted above. Taken together, the mitigations in this section convert threats that would otherwise be silent failure modes into visible, provenance-tagged, assurance-graded artefacts---an outcome the architecture treats as the realistic target, since the alternative of eliminating the threats outright is not available.
